\section{Complete Results}
\label{appx:complete_results}

This appendix provides comprehensive experimental results for all four evaluation tasks: long-term forecasting, time series classification, imputation, and anomaly detection.

\subsection{Long-term Forecasting}
\label{appx:forecasting_results}

\begin{table}[htbp]
  \caption{Full results for the long-term forecasting task. We compare extensive competitive models under different prediction lengths. The input sequence length is set to 96 for all datasets. \emph{Avg} is averaged from all four prediction lengths, that $\{96, 192, 336, 720\} $. \boldres{Red bold}: the best results; \secondres{underlined in blue}: the second best.}\label{tab:full_long_term_forecasting_results}
  \vskip 0.05in
  \centering
  \resizebox{1.0\columnwidth}{!}{
  \begin{threeparttable}
  \begin{small}
  \renewcommand{\multirowsetup}{\centering}
  \setlength{\tabcolsep}{1.2pt}
  \begin{tabular}{c|c|cc|cc|cc|cc|cc|cc|cc|cc|cc|cc|cc|cc|cc|cc}
    \toprule
    \multicolumn{2}{c}{\multirow{2}{*}{Models}} & \multicolumn{2}{c}{\rotatebox{0}{{\textbf{MAS4TS}}}} &
    \multicolumn{2}{c}{\rotatebox{0}{{MAFS}}} &
    \multicolumn{2}{c}{\rotatebox{0}{{TimeVLM}}} &
    \multicolumn{2}{c}{\rotatebox{0}{{UniTime}}} &
    \multicolumn{2}{c}{\rotatebox{0}{{iTransformer}}} &
    \multicolumn{2}{c}{\rotatebox{0}{{PatchTST}}} &
    \multicolumn{2}{c}{\rotatebox{0}{{Crossformer}}} &
    \multicolumn{2}{c}{\rotatebox{0}{{TiDE}}} &
    \multicolumn{2}{c}{\rotatebox{0}{{TimesNet}}} & 
    \multicolumn{2}{c}{\rotatebox{0}{{DLinear}}} &
    \multicolumn{2}{c}{\rotatebox{0}{{SCINet}}}&
    \multicolumn{2}{c}{\rotatebox{0}{{FEDformer}}} & 
    \multicolumn{2}{c}{\rotatebox{0}{{Stationary}}} & 
    \multicolumn{2}{c}{\rotatebox{0}{{Autoformer}}}  \\
    \multicolumn{2}{c}{} & \multicolumn{2}{c}{{(\textbf{Ours})}} &
    \multicolumn{2}{c}{{\citeyear{huang2025many}}} &
    \multicolumn{2}{c}{{\citeyear{zhong2025time}}} &
    \multicolumn{2}{c}{{\citeyear{liu2024unitime}}} &
    \multicolumn{2}{c}{{\citeyear{liu2023itransformer}}} &
    \multicolumn{2}{c}{{\citeyear{nie2023patchtst}}}&
    \multicolumn{2}{c}{{\citeyear{zhang2023crossformer}}}&
    \multicolumn{2}{c}{{\citeyear{das2023long}}}&
    \multicolumn{2}{c}{{\citeyear{wu2022timesnet}}}&
    \multicolumn{2}{c}{{\citeyear{zeng2023dlinear}}}&
    \multicolumn{2}{c}{{\citeyear{liu2022scinet}}}&
    \multicolumn{2}{c}{{\citeyear{zhou2022fedformer}}}&
    \multicolumn{2}{c}{{\citeyear{liu2022non}}}&
    \multicolumn{2}{c}{{\citeyear{wu2021autoformer}}}
    \\
    \hline    
    % \cmidrule(lr){3-4} \cmidrule(lr){5-6}\cmidrule(lr){7-8} \cmidrule(lr){9-10}\cmidrule(lr){11-12}\cmidrule(lr){13-14}\cmidrule(lr){15-16}\cmidrule(lr){17-18}\cmidrule(lr){19-20} \cmidrule(lr){21-22} \cmidrule(lr){23-24} \cmidrule(lr){25-26}
    \multicolumn{2}{c|}{Metric} & {MSE} & {MAE} & {MSE} & 
    {MAE} & {MSE} & 
    {MAE} & {MSE} & 
    {MAE} & {MSE} & 
    {MAE} & {MSE} & {MAE} & {MSE} & {MAE} & {MSE} & {MAE} & {MSE} & {MAE} & {MSE} & {MAE} & {MSE} & {MAE} & {MSE} & {MAE} & {MSE} & {MAE}& {MSE} & {MAE}\\
    \toprule

 \multirow{5}{*}{\rotatebox{90}{{Weather}}} 
    & {96} &\boldres{{0.147}}&\boldres{{0.197}}& \secondres{0.166} & \secondres{0.208} & 0.181 & 0.220 & 0.171 & 0.214 & {0.174} & 0.214&{{0.186}}&{{0.227}}&{0.195}&{0.271}& {0.202} & {0.261} & 0.172 &{{0.220}}&{0.195}&{0.252} & {0.221} &{0.306}& {0.217} &{0.296}& {0.173} &{0.223} & {0.266} &{0.336}  \\
    & {192} &\boldres{{0.192}}&\boldres{{0.244}} & 0.218 & \secondres{0.253} & 0.227 & 0.260 & 0.217 & 0.254 & {0.221} & 0.254 & {0.234}&{{0.265}}&\secondres{{0.209}}&{0.277}& {0.242} & {0.298} &{{0.219}} &{{0.261}}&{0.237}&{0.295} & {0.261} &{0.340}& {0.276} &{0.336} & {0.245} &{0.285} & {0.307} &{0.367} \\
    & {336} &\boldres{{0.245}}&\boldres{{0.286}}& 0.274 & 0.295 & 0.281 & 0.298 & 0.274 & \secondres{0.293} & {0.278} & 0.296&{{0.284}}&{0.301}&\secondres{{0.273}}&{0.332}& {0.287} & {0.335} &{{0.280}} &{{0.306}}&{0.282}&{0.331}  & {0.309} &{0.378}& {0.339} &{0.380} & {0.321} &{0.338} & {0.359} &{0.395}\\
    & {720} &\boldres{{0.328}}&\boldres{{0.339}}&0.351 & 0.345 & 0.358 & 0.348 & 0.351 & \secondres{0.343} & {0.358} & 0.347&{0.356}&{0.349}&{0.379}&{0.401}& {{0.351}} & {0.386} &{0.365} &{{0.359}}&\secondres{{0.345}}&{0.382} & {0.377} &{0.427} & {0.403} &{0.428} & {0.414} &{0.410} & {0.419} &{0.428} 
    \\
    \cmidrule(lr){2-30}
    & {Avg} &\boldres{{0.228}}&\boldres{{0.267}}& \secondres{0.252} & \secondres{0.275} & 0.262 & 0.281 & 0.253 & 0.276 & 0.258 & 0.278&{{0.265}}&{0.285}&{0.264}&{0.320}& {0.271} & {0.320} &{{0.259}} &{{0.287}}&{0.265}&{0.315}&{0.292} &{0.363} &{0.309} &{0.360} &{0.288} &{0.314} &{0.338} &{0.382}  \\
    \midrule
    
    \multirow{5}{*}{\rotatebox{90}{{Solar-Energy}}} 
    & {96} &\boldres{{0.192}}& 0.259 & \secondres{0.198} & \boldres{0.237} & 0.220 & 0.269 & 0.239 & 0.301 &0.203 &\secondres{0.237}&{0.265}&{0.323}&{{0.232}}&{0.302} &{0.312} &{0.399}&{0.373}&{0.358}&{0.290}&{0.378}&{0.237}&{0.344}&{0.286}&{0.341}&{0.321}&{0.380}&{0.456}&{0.446}
    \\
    & {192} &\boldres{{0.207}}&0.271& \secondres{0.231} & \secondres{0.263} & 0.257 & 0.293 & 0.269 & 0.316 &0.233 &\boldres{{0.261}}&{0.288}&{0.332}&{0.371}&{0.410} &{0.339} &{0.416}&{0.397}&{0.376}&{0.320}&{0.398}&{0.280}&{0.380}&{0.291}&{0.337}&{0.346}&{0.369}&{0.588}&{0.561}
    \\
    & {336} &\boldres{{0.219}}&0.283& 0.250 & \secondres{0.278} & 0.271 & 0.303 & 0.288 & 0.324 &\secondres{{0.248}} &\boldres{{0.273}}&{0.301}&{0.339}&{0.495}&{0.515}&{0.368} &{0.430}&{0.420}&{0.380}&{0.353}&{0.415}&{0.304}&{0.389}&{0.354}&{0.416}&{0.357}&{0.387}&{0.595}&{0.588}\\
    & {720} &\boldres{{0.231}}&\secondres{{0.293}}& 0.255 & 0.281 & 0.287 & 0.307 & 0.285 & 0.320 & \secondres{{0.249}} &\boldres{{0.275}}&{0.295}&{0.336}&{0.526}&{0.542}&{0.370} &{0.425}&{0.420}&{0.381}&{0.357}&{0.413}&{0.308}&{0.388}
    &{0.380}&{0.437}&{0.375}&{0.424}&{0.733}&{0.633}
    \\
    \cmidrule(lr){2-30}
    & {Avg}&\boldres{{0.212}} &0.276 & 0.234 & \secondres{0.265} & 0.259 & 0.293 & 0.270 & 0.315 &\secondres{{0.233}} &\boldres{{0.262}}&{0.287}&{0.333}&{0.406}&{0.442}&{0.347} &{0.417}&{0.403}&{0.374}&{0.330}&{0.401}&{0.282} &{0.375}&{0.328} &{0.383} &{0.350} &{0.390} &{0.586} &{0.557} \\
    \midrule

    \multirow{5}{*}{\rotatebox{90}{{Exchange}}} 
    & {96}  &0.088&0.208& \secondres{0.084}	& \secondres{0.204} & \boldres{0.082} & \boldres{0.198} & 0.125 & 0.245& 0.086 & 0.206 & 0.088 &0.205 & {0.256} & {0.367} & {0.094} & {0.218} & {0.107} & {0.234} & {0.088} & {0.218} & {0.267} & {0.396} & {0.148} & {0.278} & {0.111} & {0.237} & {0.197} & {0.323} 
    \\
    & {192}  &0.181&{{0.303}} & 0.180 & 0.304 & \boldres{0.175} & \boldres{0.296} & 0.236 & 0.343 & {0.177} & \secondres{{0.299}} & \secondres{0.176} & \secondres{{0.299}} & {0.470} & {0.509} & {0.184} & {0.307} & {0.226} & {0.344} & 0.176 & {0.315} & {0.351} & {0.459} & {0.271} & {0.315} & {0.219} & {0.335} & {0.300} & {0.369} 
    \\
    & {336}  &\boldres{{0.293}}&\boldres{{0.396}} & 0.316 & 0.406 & 0.324 & 0.412 & 0.420 & 0.470 & {0.331} & \secondres{{0.417}}& \secondres{{0.301}} & \secondres{{0.397}} & {1.268} & {0.883} & {0.349} & {0.431} & {0.367} & {0.448} & {0.313} & {0.427} & {1.324} & {0.853} & {0.460} & {0.427} & {0.421} & {0.476} & {0.509} & {0.524}
    \\
    & {720}  &\boldres{{0.810}}&\boldres{{0.682}} & 1.002 & 0.755 & 0.840 & \secondres{0.689} & 1.078 & 0.791 &{{0.847}} & 0.691 & {0.901} & {0.714} & {1.767} & {1.068} & {0.852} & {0.698} & {0.964} & {0.746} & \secondres{{0.839}} & {{0.695}} & {1.058} & {0.797} & {1.195} & {{0.695}} & {1.092} & {0.769} & {1.447} & {0.941}
    \\
    \cmidrule(lr){2-30}
    & {Avg} & \boldres{{0.343}} & \boldres{0.397} & 0.396 & 0.417 & 0.355 & \secondres{0.399} & 0.465 & 0.462 &{{0.360}} & 0.403 & {0.367} & 0.404 & {0.940} & {0.707} & {0.370} & {0.413} & {0.416} & {0.443} & \secondres{{0.354}} & {0.414} & {0.750} & {0.626} & {0.519} & {0.429} & {0.461} & {0.454} & {0.613} & {0.539} \\
    \midrule

    \multirow{5}{*}{\rotatebox{90}{{ETTh1}}} 
    & {96} &\secondres{{0.378}}&\secondres{{0.399}}& 0.389 & 0.403 & \boldres{0.376} & \boldres{0.395} & 0.397 & 0.418 & 0.386 & \secondres{{0.405}}&{{0.460}}&{{0.447}}&{0.423}&{0.448} & {0.479}& {0.464}  &{{0.384}} &{{0.402}}&{0.397}&{0.412}&{0.654} &{0.599} 
     &{0.395} &{0.424} &{0.513} &{0.491} &{0.449} &{0.459} \\
    & {192} &\boldres{{0.419}}&\secondres{{0.428}}& 0.450 & 0.438 & \secondres{0.426} & \boldres{0.424} & 0.434 & 0.439 & {0.441} & {0.512}&{{0.477}}&0.429&{0.471}&{0.474}& {0.525} & {0.492} &0.436 &{{0.429}}&{0.446}&{0.441} &{0.719} &{0.631}
     &{0.469} &{0.470} &{0.534} &{0.504} &{0.500} &{0.482} 
    \\
    & {336} &\boldres{{0.459}}&\boldres{0.445} & 0.475 & \secondres{0.447} & 0.469 & 0.449 & \secondres{0.468} & 0.457 & 0.487 & 0.458&{0.546}&{0.496}&{0.570}&{0.546}& {0.565} & {0.515} &{0.491} &{0.469}&{{0.489}}&{{0.467}}&{0.778} &{0.659}
    &{0.530} &{0.499} &{0.588} &{0.535} &{0.521} &{0.496} 
    \\
    & {720} & 0.503 & {{0.501}} & \secondres{0.487}	& \boldres{0.472} & 0.559 & 0.504 & \boldres{0.469} & \secondres{0.477} & 0.503 & 0.491&{0.544}&{{0.517}}&{0.653}&{0.621}& {0.594} & {0.558} &{0.521} &0.500&{{0.513}}&{0.510}&{0.836} &{0.699}
    &{0.598} &{0.544} &{0.643} &{0.616} &{{0.514}} &{0.512}  
    \\
    \cmidrule(lr){2-30}
    & {Avg} &\boldres{{0.440}}&\boldres{{0.438}}& 0.450 & \secondres{0.440} & 0.457 & 0.443 & 0.442 & 0.448 & 0.454 & 0.447&{0.516}&{{0.484}}&{0.529}&{0.522}& {0.541} & {0.507} &{0.458} &{{0.450}}&{{0.461}}&{0.457}&{0.747} &{0.647}&{0.498} &{0.484} &{0.570} &{0.537} &{0.496} &{0.487} \\
    \midrule

    \multirow{5}{*}{\rotatebox{90}{{ETTh2}}} 
    & {96}&\secondres{{0.296}}&\secondres{{0.345}}& 0.298 & 0.346 & \boldres{0.290} & \boldres{0.339} & 0.296 & 0.345 & 0.297 & 0.349 &{0.308}&{0.355}&{0.745}&{0.584}&{0.400} & {0.440}  & {{0.340}} & {{0.374}}&{0.340}&{0.394}&{0.707} &{0.621}  &{0.358} &{0.397} &{0.476} &{0.458} &{0.346} &{0.388}  \\
    & {192} &0.386&0.399& 0.380 & 0.398 & \boldres{0.374} & \boldres{0.389} & \secondres{0.374} & \secondres{0.394} & 0.380 & 0.400&{0.393}&{0.405}&{0.877}&{0.656}& {0.528} & {0.509} & {{0.402}} & {{0.414}}&{0.482}&{0.479}&{0.860} &{0.689} &{0.429} &{0.439} &{0.512} &{0.493} &{0.456} &{0.452}  \\
    & {336} &\boldres{{0.407}}&\boldres{{0.422}} & \secondres{0.413} & \secondres{0.427} & 0.433 & 0.440 & 0.415 & 0.427 & {{0.428}} & 0.432&0.427 &{0.436}&{1.043}&{0.731}& {0.643} & {0.571}  & {{0.452}} & {{0.452}}&{0.591}&{0.541} &{1.000} &{0.744}&{0.496} &{0.487} &{0.552} &{0.551} &{0.482} &{0.486} \\
    & {720} &\secondres{{0.427}}&\boldres{{0.442}}& 0.428 & \secondres{0.444} & 0.438 & 0.452 & \boldres{0.425} & 0.444 & 0.427 & 0.445&{0.436}&{0.450}&{1.104}&{0.763}& {0.874} & {0.679} & {{0.462}} & {{0.468}}&{0.839}&{0.661} &{1.249} &{0.838}&{0.463} &{0.474}&{0.562} &{0.560} &{0.515} &{0.511}  
    \\
    \cmidrule(lr){2-30}
    & {Avg} &\boldres{{0.379}}&\boldres{{0.402}} & 0.380 & \secondres{0.403} & 0.384 & 0.405 & \secondres{0.379} & 0.403 & 0.383 & 0.407&{0.391}&{0.411}&{0.942}&{0.684}& {0.611} & {0.550}  &{{0.414}} &{{0.427}}&{0.563}&{0.519} &{0.954} &{0.723}&{0.437} &{0.449} &{0.526} &{0.516} &{0.450} &{0.459} \\
    \midrule

    \multirow{5}{*}{\rotatebox{90}{{ETTm1}}} 
    & {96} &\boldres{{0.296}}&\boldres{{0.341}}& 0.329 & \secondres{0.362} & \secondres{0.326} & 0.365 & 0.322 & 0.363 & 0.334 & 0.368&{{0.352}}&{0.374}&{0.404}&{0.426}& {0.364} & {0.387} &{{0.338}} &{{0.375}}&{0.346}&{0.374} &{0.418} &{0.438} &{0.379} &{0.419} &{0.386} &{0.398} &{0.505} &{0.475} 
    \\
    & {192} &\boldres{{0.341}}&\boldres{{0.366}}&0.370 & 0.385 & 0.368 & \secondres{0.384} & \secondres{0.366} & 0.387 & {0.390} & {0.393}&0.374&0.387&{0.450}&{0.451}&{0.398} & {0.404} &{{0.374}} &{{0.387}}&{0.382}&{0.391} &{0.439} &{0.450}&{0.426} &{0.441} &{0.459} &{0.444} &{0.553} &{0.496} 
    \\
    & {336} &\boldres{{0.385}}&\boldres{{0.390}}& 0.409 & 0.410 & \secondres{0.394} & \secondres{0.405} & 0.398 & 0.407 & {0.426} & {0.420}&{{0.421}}&{0.414}&{0.532}&{0.515}& {0.428} & {0.425} &0.410 &0.411&{0.415}&{0.415} &{0.490} &{0.485}&{0.445} &{0.459} &{0.495} &{0.464} &{0.621} &{0.537}  
    \\
    & {720} &\boldres{{0.445}}&\boldres{{0.426}}&0.478 & 0.451 & 0.462 & \secondres{0.440} & \secondres{0.454} & 0.440 & {0.491} & {0.459}&0.462&0.449&{0.666}&{0.589}& {0.487} & {0.461} &{{0.478}} &{0.450}&{0.473}&{0.451} &{0.595} &{0.550}&{0.543} &{0.490} &{0.585} &{0.516} &{0.671} &{0.561} 
    \\
    \cmidrule(lr){2-30}
    & {Avg} &\boldres{{0.367}}&\boldres{{0.381}}& 0.397 & 0.402 & 0.388 & \secondres{0.398} & \secondres{0.385} & 0.399 & {0.407} & {0.410}&{{0.406}}&{0.407}&{0.513}&{0.495}& {0.419} & {0.419} &0.400 &0.406&{0.404}&{0.408} &{0.485} &{0.481} &{0.448} &{0.452} &{0.481} &{0.456} &{0.588} &{0.517} \\
    \midrule

    \multirow{5}{*}{\rotatebox{90}{{ETTm2}}} 
    & {96} &\boldres{{0.175}}&\boldres{{0.260}}& 0.180 & 0.264 & \secondres{0.176} & \secondres{0.261} & 0.183 & 0.266 & 0.180 & 0.264 &{0.183} &{0.270}&{0.287}&{0.366}& {0.207} & {0.305} &{{0.187}} &{0.267}&{0.193}&{0.293}&{0.286} &{0.377}
    &{0.203} &{0.287} &{0.192} &{0.274} &{0.255} &{0.339} 
    \\
    & {192} &\boldres{{0.237}}&\boldres{{0.302}}&0.245 & 0.304 & \secondres{0.240} & \secondres{0.302} & 0.251 & 0.310 & {0.250} & 0.309&{0.255}&{0.314}&{0.414}&{0.492}& {0.290} & {0.364} &0.249 &{{0.309}}&{0.284}&{0.361}&{0.399} &{0.445}
    &{0.269} &{0.328} &{0.280} &{0.339} &{0.281} &{0.340} 
    \\
    & {336} &\secondres{0.303}&\secondres{0.343}& 0.316	& 0.351 & \boldres{0.299} & \boldres{0.339} & 0.319 & 0.351 & {{0.311}} & 0.348 &0.309 & 0.347&{0.597}&{0.542} & {0.377} & {0.422} &{{0.321}} &{{0.351}}&{0.382}&{0.429}&{0.637} &{0.591}
    &{0.325} &{0.366} &{0.334} &{0.361} &{0.339} &{0.372} 
    \\
    & {720} &\boldres{{0.396}}& \boldres{0.400} & 0.409 & \secondres{0.403} & \secondres{0.400} & {0.416} & 0.420 & 0.410 & 0.412 & 0.407&{{0.412}}&{0.404}&{1.730}&{1.042}& {0.558} & {0.524} &{{0.408}} &0.403&{0.558}&{0.525}&{0.960} &{0.735}
     &{0.421} &{0.415} &{0.417} &{0.413} &{0.433} &{0.432} \\
    \cmidrule(lr){2-30}
    &{Avg} &\boldres{{0.278}}&\boldres{0.326} & 0.288 & 0.330 & \secondres{0.279} & \secondres{0.329} & 0.293 & 0.334 & 0.288 & \secondres{{0.332}}&{0.290}&{0.334}&{0.757}&{0.610}& {0.358} & {0.404} &{{0.291}} &{{0.333}}&{0.354}&{0.402} &{0.954} &{0.723}&{0.305} &{0.349} &{0.306} &{0.347} &{0.327} &{0.371} \\
%    \bottomrule
\midrule
\multicolumn{2}{c|}{{$1^{st}$ Count}} & {27} & {22} & 0 & 2 & 6 & 7 & 2 & 0 & {0} & {4} & {0} & {0} & {0} & {0} & {0} & {0} & {0} & {0} & {0} & {0} & {0} & {0} & {0} & {0} & {0} & {0} & {0} & {0} \\
\bottomrule
  \end{tabular}
    \end{small}
  \end{threeparttable}
  }
\end{table}

Table~\ref{tab:full_long_term_forecasting_results} presents complete forecasting results across all datasets and prediction horizons $H \in \{96, 192, 336, 720\}$. The input sequence length is fixed at $L = 96$ for all methods following prior work~\citep{wu2022timesnet}. MAS4TS consistently achieves competitive or superior performance compared to both task-specific models (e.g., PatchTST, iTransformer) and pretrained LM-based methods (e.g., UniTime, TimeVLM).

\subsection{Classification}
\label{appx:classification_results}
\begin{table}[tbp]
  \caption{Full results for the classification task. $\ast.$ in the model names indicates the name of $\ast$former. We report the classification accuracy (\%) as the result. The standard deviation is within 0.1\%. {Red Bold}: the best results; {underlined in blue}: the second best.}
  \label{tab:full_classification_results}
  \vskip 0.05in
  \centering
  \resizebox{\columnwidth}{!}{
  \begin{threeparttable}
  \begin{small}
  \renewcommand{\multirowsetup}{\centering}
  \setlength{\tabcolsep}{2pt}
  \begin{tabular}{c|cccccccccccccccccc}
    \toprule
    \scalebox{0.8}{Methods} & \scalebox{0.7}{DTW} & \scalebox{0.6}{XGBoost} & \scalebox{0.6}{LSTM} & \scalebox{0.6}{LSTNet} & \scalebox{0.6}{LSSL} & \scalebox{0.6}{TCN} & \scalebox{0.7}{Trans.} & \scalebox{0.7}{Re.} & \scalebox{0.7}{In.} & \scalebox{0.7}{Pyra.} & \scalebox{0.7}{Auto.} &  \scalebox{0.7}{FED.} & \scalebox{0.7}{ETS.} &  \scalebox{0.7}{iTrans.}& \scalebox{0.6}{DLinear} & \scalebox{0.6}{LightTS}&  \scalebox{0.7}{TiDE} &  \scalebox{0.7}{\textbf{MAS4TS}} \\
	\scalebox{0.8}{Datasets} & \scalebox{0.6}{\citeyearpar{Berndt1994UsingDT}} & \scalebox{0.6}{\citeyearpar{Chen2016XGBoostAS}} & \scalebox{0.6}{\citeyearpar{Hochreiter1997LongSM}} & 
	\scalebox{0.6}{\citeyearpar{2018Modeling}} & 
	\scalebox{0.6}{\citeyearpar{gu2022efficiently}} & 
	\scalebox{0.6}{\citeyearpar{Franceschi2019UnsupervisedSR}} & \scalebox{0.6}{\citeyearpar{vaswani2017attention}} & 
	\scalebox{0.6}{\citeyearpar{kitaev2020reformer}} & \scalebox{0.6}{\citeyearpar{zhou2021informer}} & \scalebox{0.6}{\citeyearpar{liu2022pyraformer}} &
	\scalebox{0.6}{\citeyearpar{wu2021autoformer}} & 
	\scalebox{0.6}{\citeyearpar{zhou2022fedformer}} & \scalebox{0.6}{\citeyearpar{woo2022etsformer}} & \scalebox{0.6}{\citeyearpar{liu2023itransformer}} & 
	\scalebox{0.6}{\citeyearpar{zeng2023dlinear}} & \scalebox{0.6}{\citeyearpar{campos2023lightts}} & \scalebox{0.6}{\citeyearpar{das2023long}}& \scalebox{0.6}{\textbf{(Ours)}} \\
    \midrule
	\scalebox{0.7}{EthanolConcentration} & \scalebox{0.8}{32.3} & {\scalebox{0.8}{43.7}} & \scalebox{0.8}{32.3} & \scalebox{0.8}{39.9} & \scalebox{0.8}{31.1}&  \scalebox{0.8}{28.9} & \scalebox{0.8}{32.7} &\scalebox{0.8}{31.9} &\scalebox{0.8}{31.6}   &\scalebox{0.8}{30.8} &\scalebox{0.8}{31.6} & \scalebox{0.8}{28.1}&\scalebox{0.8}{31.2}  & \scalebox{0.8}{28.1}& \scalebox{0.8}{32.6} &\scalebox{0.8}{29.7} & \scalebox{0.8}{27.1}& {\scalebox{0.8}{33.46}}\\
	\scalebox{0.7}{FaceDetection} & \scalebox{0.8}{52.9} & \scalebox{0.8}{63.3} & \scalebox{0.8}{57.7} & \scalebox{0.8}{65.7} & \scalebox{0.8}{66.7} & \scalebox{0.8}{52.8} & \scalebox{0.8}{67.3} & {\scalebox{0.8}{68.6}} &\scalebox{0.8}{67.0} &\scalebox{0.8}{65.7} &\scalebox{0.8}{68.4} &\scalebox{0.8}{66.0} & \scalebox{0.8}{66.3} & \scalebox{0.8}{66.3}&\scalebox{0.8}{68.0} &\scalebox{0.8}{67.5} & \scalebox{0.8}{65.3}& {\scalebox{0.8}{69.38}}  \\
	\scalebox{0.7}{Handwriting} & \scalebox{0.8}{28.6} & \scalebox{0.8}{15.8} & \scalebox{0.8}{15.2} & \scalebox{0.8}{25.8} & \scalebox{0.8}{24.6} & {\scalebox{0.8}{53.3}} & \scalebox{0.8}{32.0} & \scalebox{0.8}{27.4} &\scalebox{0.8}{32.8} &\scalebox{0.8}{29.4} &\scalebox{0.8}{36.7} &\scalebox{0.8}{28.0} &  \scalebox{0.8}{32.5} & \scalebox{0.8}{24.2}& \scalebox{0.8}{27.0} &\scalebox{0.8}{26.1} & \scalebox{0.8}{23.2}& {\scalebox{0.8}{33.64}} \\
	\scalebox{0.7}{Heartbeat} & \scalebox{0.8}{71.7}  & \scalebox{0.8}{73.2} & \scalebox{0.8}{72.2} & \scalebox{0.8}{77.1} & \scalebox{0.8}{72.7}& \scalebox{0.8}{75.6} & \scalebox{0.8}{76.1} & \scalebox{0.8}{77.1} &{\scalebox{0.8}{80.5}} &\scalebox{0.8}{75.6} &\scalebox{0.8}{74.6} &\scalebox{0.8}{73.7} &  \scalebox{0.8}{71.2} & \scalebox{0.8}{75.6}& \scalebox{0.8}{75.1} &\scalebox{0.8}{75.1} & \scalebox{0.8}{74.6}& {\scalebox{0.8}{78.04}}    \\
	\scalebox{0.7}{JapaneseVowels} & \scalebox{0.8}{94.9} & \scalebox{0.8}{86.5} & \scalebox{0.8}{79.7} & \scalebox{0.8}{98.1} & \scalebox{0.8}{98.4} & {\scalebox{0.8}{98.9}} & \scalebox{0.8}{98.7} & \scalebox{0.8}{97.8} &{\scalebox{0.8}{98.9}} &\scalebox{0.8}{98.4} &\scalebox{0.8}{96.2} &\scalebox{0.8}{98.4} & \scalebox{0.8}{95.9} & \scalebox{0.8}{96.6}& \scalebox{0.8}{96.2} &\scalebox{0.8}{96.2} & \scalebox{0.8}{95.6}& {\scalebox{0.8}{97.56}}  \\
	\scalebox{0.7}{SelfRegulationSCP1} & \scalebox{0.8}{77.7}  & \scalebox{0.8}{84.6} & \scalebox{0.8}{68.9} & \scalebox{0.8}{84.0} & \scalebox{0.8}{90.8} & \scalebox{0.8}{84.6} & {\scalebox{0.8}{92.2}} & \scalebox{0.8}{90.4} &\scalebox{0.8}{90.1} &\scalebox{0.8}{88.1} &\scalebox{0.8}{84.0} &\scalebox{0.8}{88.7} & \scalebox{0.8}{89.6} & \scalebox{0.8}{90.2}& \scalebox{0.8}{87.3} &\scalebox{0.8}{89.8} & \scalebox{0.8}{89.2}& {\scalebox{0.8}{91.81}} \\
    \scalebox{0.7}{SelfRegulationSCP2} & \scalebox{0.8}{53.9} & \scalebox{0.8}{48.9} & \scalebox{0.8}{46.6} & \scalebox{0.8}{52.8} & \scalebox{0.8}{52.2} & \scalebox{0.8}{55.6} & \scalebox{0.8}{53.9} & {\scalebox{0.8}{56.7}} &\scalebox{0.8}{53.3} &\scalebox{0.8}{53.3} &\scalebox{0.8}{50.6} &\scalebox{0.8}{54.4} & \scalebox{0.8}{55.0} & \scalebox{0.8}{54.4}& \scalebox{0.8}{50.5} &\scalebox{0.8}{51.1} & \scalebox{0.8}{53.4}& {\scalebox{0.8}{55.55}}  \\
    \scalebox{0.7}{UWaveGestureLibrary} & {\scalebox{0.8}{90.3}} & \scalebox{0.8}{75.9} & \scalebox{0.8}{41.2} & \scalebox{0.8}{87.8} & \scalebox{0.8}{85.9} & \scalebox{0.8}{88.4} & \scalebox{0.8}{85.6} & \scalebox{0.8}{85.6} &\scalebox{0.8}{85.6} &\scalebox{0.8}{83.4} &\scalebox{0.8}{85.9} &\scalebox{0.8}{85.3} & \scalebox{0.8}{85.0} & \scalebox{0.8}{85.9}& \scalebox{0.8}{82.1} &\scalebox{0.8}{80.3} & \scalebox{0.8}{84.9}& {\scalebox{0.8}{86.56}} \\
    \midrule
    \scalebox{0.8}{Average Accuracy} & \scalebox{0.8}{62.78} & \scalebox{0.8}{61.48} & \scalebox{0.8}{51.72} & \scalebox{0.8}{66.40} & \scalebox{0.8}{65.30} & \scalebox{0.8}{67.26} & \scalebox{0.8}{67.31} & {\scalebox{0.8}{66.93}} &\scalebox{0.8}{59.93} &\scalebox{0.8}{65.58} &\scalebox{0.8}{66.00} &\scalebox{0.8}{65.32} & \scalebox{0.8}{65.83} & \scalebox{0.8}{65.16}& \scalebox{0.8}{64.85} &\scalebox{0.8}{64.47} & \scalebox{0.8}{64.16}& {\scalebox{0.8}{\textbf{68.25}}} \\
	\bottomrule
  \end{tabular}
    \end{small}
  \end{threeparttable}
  }
\end{table}

Table~\ref{tab:full_classification_results} shows classification accuracy on 8 representative datasets from the UEA multivariate time series archive~\citep{bagnall2018uea}. These datasets span diverse application domains including chemical sensing, video analysis, motion capture, medical diagnostics, and audio recognition.

\subsection{Imputation}
\label{appx:imputation_results}
\begin{table}[!ht]
\centering
\caption{Full results for the imputation task. We randomly mask 12.5\%, 25\%, 37.5\% and 50\% time points to compare the model performance under different missing degrees. $*$ in the Transformers indicates the name of $*$former.}
\label{tab:full_imputation_results}
\resizebox{\textwidth}{!}{
\begin{tabular}{cc|cc|cc|cc|cc|cc|cc|cc|cc|cc|cc|cc}
\hline
\multicolumn{2}{c|}{Models} & \multicolumn{2}{c|}{MAS4TS} & \multicolumn{2}{c|}{ETS.} & \multicolumn{2}{c|}{LightTS$^*$} & \multicolumn{2}{c|}{DLinear$^*$} & \multicolumn{2}{c|}{FED.} & \multicolumn{2}{c|}{Pyra.} & \multicolumn{2}{c|}{In.} & \multicolumn{2}{c|}{Re.} & \multicolumn{2}{c|}{LSTM} & \multicolumn{2}{c|}{TCN} & \multicolumn{2}{c}{LSSL} \\
\multicolumn{2}{c|}{} & \multicolumn{2}{c|}{(Ours)} & \multicolumn{2}{c|}{(2022)} & \multicolumn{2}{c|}{(2022)} & \multicolumn{2}{c|}{(2023)} & \multicolumn{2}{c|}{(2022)} & \multicolumn{2}{c|}{(2021a)} & \multicolumn{2}{c|}{(2021)} & \multicolumn{2}{c|}{(2020)} & \multicolumn{2}{c|}{(1997)} & \multicolumn{2}{c|}{(2019)} & \multicolumn{2}{c}{(2022)} \\
\hline
\multicolumn{2}{c|}{Mask Ratio} & MSE & MAE & MSE & MAE & MSE & MAE & MSE & MAE & MSE & MAE & MSE & MAE & MSE & MAE & MSE & MAE & MSE & MAE & MSE & MAE & MSE & MAE \\
\hline
\multirow{5}{*}{\rotatebox{90}{ETTm1}} & 12.5\% & 0.046 & 0.137 & 0.067 & 0.188 & 0.075 & 0.180 & 0.058 & 0.162 & \secondres{0.035} & \secondres{0.135} & 0.670 & 0.541 & 0.047 & 0.155 & \boldres{0.032} & \boldres{0.126} & 0.974 & 0.780 & 0.510 & 0.493 & 0.101 & 0.231 \\
& 25\% & \secondres{0.049} & \boldres{0.143} & 0.096 & 0.229 & 0.093 & 0.206 & 0.080 & 0.193 & 0.052 & 0.166 & 0.689 & 0.553 & 0.063 & 0.180 & \boldres{0.042} & \secondres{0.146} & 1.032 & 0.807 & 0.518 & 0.500 & 0.106 & 0.235 \\
& 37.5\% & \secondres{0.054} & \boldres{0.150} & 0.133 & 0.271 & 0.113 & 0.231 & 0.103 & 0.219 & 0.069 & 0.191 & 0.737 & 0.581 & 0.079 & 0.200 & \boldres{0.052} & \secondres{0.158} & 0.999 & 0.792 & 0.516 & 0.499 & 0.116 & 0.246 \\
& 50\% & \secondres{0.061} & \boldres{0.159} & 0.186 & 0.323 & 0.134 & 0.255 & 0.132 & 0.248 & 0.089 & 0.218 & 0.770 & 0.605 & 0.093 & 0.218 & \boldres{0.063} & \secondres{0.173} & 0.952 & 0.763 & 0.519 & 0.496 & 0.129 & 0.260 \\
& Avg & \secondres{0.052} & \boldres{0.147} & 0.120 & 0.253 & 0.104 & 0.218 & 0.093 & 0.206 & 0.062 & 0.177 & 0.717 & 0.570 & 0.071 & 0.188 & \boldres{0.050} & \secondres{0.154} & 0.989 & 0.786 & 0.516 & 0.497 & 0.113 & 0.254 \\
\hline
\multirow{5}{*}{\rotatebox{90}{ETTm2}} & 12.5\% & \boldres{0.026} & \boldres{0.095} & 0.108 & 0.239 & \secondres{0.034} & \secondres{0.127} & 0.062 & 0.166 & 0.056 & 0.159 & 0.394 & 0.470 & 0.133 & 0.270 & 0.108 & 0.228 & 1.013 & 0.805 & 0.307 & 0.441 & 0.150 & 0.298 \\
& 25\% & \boldres{0.029} & \boldres{0.100} & 0.164 & 0.294 & \secondres{0.042} & \secondres{0.143} & 0.085 & 0.196 & 0.080 & 0.195 & 0.421 & 0.482 & 0.135 & 0.272 & 0.136 & 0.262 & 1.039 & 0.814 & 0.263 & 0.402 & 0.159 & 0.306 \\
& 37.5\% & \boldres{0.031} & \boldres{0.105} & 0.237 & 0.356 & \secondres{0.051} & \secondres{0.159} & 0.106 & 0.222 & 0.110 & 0.231 & 0.478 & 0.521 & 0.155 & 0.293 & 0.175 & 0.300 & 0.917 & 0.744 & 0.250 & 0.396 & 0.180 & 0.321 \\
& 50\% & \boldres{0.034} & \boldres{0.111} & 0.323 & 0.421 & \secondres{0.059} & \secondres{0.174} & 0.131 & 0.247 & 0.156 & 0.276 & 0.568 & 0.560 & 0.200 & 0.333 & 0.211 & 0.329 & 1.140 & 0.835 & 0.246 & 0.389 & 0.210 & 0.353 \\
& Avg & \boldres{0.030} & \boldres{0.103} & 0.208 & 0.327 & \secondres{0.046} & \secondres{0.151} & 0.096 & 0.208 & 0.101 & 0.215 & 0.465 & 0.508 & 0.156 & 0.292 & 0.157 & 0.280 & 1.027 & 0.800 & 0.266 & 0.407 & 0.175 & 0.324 \\
\hline
\multirow{5}{*}{\rotatebox{90}{ETTh1}} & 12.5\% & 0.104 & 0.211 & 0.126 & 0.263 & 0.240 & 0.345 & 0.151 & 0.267 & \boldres{0.070} & \boldres{0.190} & 0.857 & 0.609 & 0.114 & 0.234 & \secondres{0.074} & \secondres{0.194} & 1.265 & 0.896 & 0.599 & 0.554 & 0.422 & 0.461 \\
& 25\% & 0.119 & \boldres{0.227} & 0.169 & 0.304 & 0.265 & 0.364 & 0.180 & 0.292 & \secondres{0.106} & 0.236 & 0.829 & 0.672 & 0.140 & 0.262 & \boldres{0.102} & \secondres{0.227} & 1.262 & 0.883 & 0.610 & 0.567 & 0.412 & 0.456 \\
& 37.5\% & 0.137 & \boldres{0.243} & 0.220 & 0.347 & 0.296 & 0.382 & 0.215 & 0.318 & \boldres{0.124} & \secondres{0.258} & 0.830 & 0.675 & 0.174 & 0.293 & \secondres{0.135} & 0.261 & 1.200 & 0.867 & 0.628 & 0.577 & 0.421 & 0.461 \\
& 50\% & \boldres{0.158} & \boldres{0.261} & 0.293 & 0.402 & 0.334 & 0.404 & 0.257 & 0.347 & \secondres{0.165} & \secondres{0.299} & 0.854 & 0.691 & 0.215 & 0.325 & 0.179 & \secondres{0.298} & 1.174 & 0.849 & 0.648 & 0.587 & 0.443 & 0.473 \\
& Avg & 0.129 & \boldres{0.235} & 0.202 & 0.329 & 0.284 & 0.373 & 0.201 & 0.306 & \boldres{0.117} & 0.246 & 0.842 & 0.682 & 0.161 & 0.279 & \secondres{0.122} & \secondres{0.245} & 1.225 & 0.873 & 0.621 & 0.571 & 0.424 & 0.481 \\
\hline
\multirow{5}{*}{\rotatebox{90}{ETTh2}} & 12.5\% & \boldres{0.062} & \boldres{0.161} & 0.187 & 0.319 & 0.101 & 0.231 & 0.100 & 0.216 & \secondres{0.095} & \secondres{0.212} & 0.976 & 0.754 & 0.305 & 0.431 & 0.163 & 0.289 & 2.060 & 1.120 & 0.410 & 0.494 & 0.521 & 0.555 \\
& 25\% & \boldres{0.063} & \boldres{0.160} & 0.279 & 0.390 & \secondres{0.115} & \secondres{0.246} & 0.127 & 0.247 & 0.137 & 0.258 & 1.037 & 0.774 & 0.322 & 0.444 & 0.206 & 0.331 & 2.007 & 1.105 & 0.419 & 0.490 & 0.487 & 0.535 \\
& 37.5\% & \boldres{0.068} & \boldres{0.168} & 0.400 & 0.465 & \secondres{0.126} & \secondres{0.257} & 0.158 & 0.276 & 0.187 & 0.304 & 1.107 & 0.800 & 0.353 & 0.462 & 0.252 & 0.370 & 2.033 & 1.111 & 0.429 & 0.498 & 0.487 & 0.529 \\
& 50\% & \boldres{0.074} & \boldres{0.176} & 0.602 & 0.572 & \secondres{0.136} & \secondres{0.268} & 0.183 & 0.299 & 0.232 & 0.341 & 1.193 & 0.838 & 0.369 & 0.472 & 0.316 & 0.419 & 2.054 & 1.119 & 0.467 & 0.529 & 0.484 & 0.523 \\
& Avg & \boldres{0.067} & \boldres{0.166} & 0.367 & 0.436 & \secondres{0.119} & \secondres{0.250} & 0.142 & 0.259 & 0.163 & 0.279 & 1.079 & 0.792 & 0.337 & 0.452 & 0.234 & 0.352 & 2.039 & 1.114 & 0.431 & 0.503 & 0.495 & 0.475 \\
\hline
\multirow{5}{*}{\rotatebox{90}{Electricity}} & 12.5\% & \boldres{0.070} & \boldres{0.188} & 0.196 & 0.321 & 0.102 & 0.229 & \secondres{0.092} & \secondres{0.214} & 0.107 & 0.237 & 0.297 & 0.383 & 0.218 & 0.326 & 0.190 & 0.308 & 0.277 & 0.366 & 0.621 & 0.620 & 0.217 & 0.341 \\
& 25\% & \boldres{0.079} & \boldres{0.199} & 0.207 & 0.332 & 0.121 & 0.252 & \secondres{0.118} & \secondres{0.247} & 0.120 & 0.251 & 0.294 & 0.380 & 0.219 & 0.326 & 0.197 & 0.312 & 0.281 & 0.369 & 0.559 & 0.585 & 0.219 & 0.341 \\
& 37.5\% & \boldres{0.090} & \boldres{0.201} & 0.219 & 0.344 & 0.141 & 0.273 & 0.144 & 0.276 & \secondres{0.136} & \secondres{0.266} & 0.296 & 0.381 & 0.222 & 0.328 & 0.203 & 0.315 & 0.275 & 0.364 & 0.567 & 0.588 & 0.223 & 0.343 \\
& 50\% & \boldres{0.107} & \boldres{0.221} & 0.235 & 0.357 & 0.160 & 0.293 & 0.175 & 0.305 & \secondres{0.158} & \secondres{0.284} & 0.299 & 0.383 & 0.228 & 0.331 & 0.210 & 0.319 & 0.273 & 0.361 & 0.581 & 0.597 & 0.229 & 0.347 \\
& Avg & \boldres{0.086} & \boldres{0.202} & 0.214 & 0.339 & 0.131 & 0.262 & 0.132 & 0.260 & \secondres{0.130} & \secondres{0.259} & 0.297 & 0.382 & 0.222 & 0.328 & 0.200 & 0.313 & 0.277 & 0.365 & 0.582 & 0.597 & 0.222 & 0.293 \\
\hline
\multirow{5}{*}{\rotatebox{90}{Weather}} & 12.5\% & \secondres{0.035} & 0.090 & 0.057 & 0.141 & 0.047 & 0.101 & 0.039 & \secondres{0.084} & 0.041 & 0.107 & 0.140 & 0.220 & 0.037 & 0.093 & \boldres{0.031} & \boldres{0.076} & 0.296 & 0.379 & 0.176 & 0.287 & 0.036 & 0.095 \\
& 25\% & \boldres{0.035} & \boldres{0.062} & 0.065 & 0.155 & 0.052 & 0.111 & 0.048 & 0.103 & 0.064 & 0.163 & 0.147 & 0.229 & 0.042 & 0.100 & \secondres{0.035} & \secondres{0.082} & 0.327 & 0.409 & 0.187 & 0.293 & 0.042 & 0.104 \\
& 37.5\% & \boldres{0.037} & \boldres{0.065} & 0.081 & 0.180 & 0.058 & 0.121 & 0.057 & 0.117 & 0.107 & 0.229 & 0.156 & 0.240 & 0.049 & 0.111 & \secondres{0.040} & \secondres{0.091} & 0.406 & 0.463 & 0.172 & 0.281 & 0.047 & 0.112 \\
& 50\% & \boldres{0.036} & \boldres{0.063} & 0.102 & 0.207 & 0.065 & 0.133 & 0.066 & 0.134 & 0.183 & 0.312 & 0.164 & 0.249 & 0.053 & 0.114 & \secondres{0.046} & \secondres{0.099} & 0.431 & 0.483 & 0.195 & 0.303 & 0.054 & 0.123 \\
& Avg & \boldres{0.036} & \boldres{0.070} & 0.076 & 0.171 & 0.055 & 0.117 & 0.052 & 0.110 & 0.099 & 0.203 & 0.152 & 0.235 & 0.045 & 0.104 & \secondres{0.038} & \secondres{0.087} & 0.365 & 0.434 & 0.183 & 0.291 & 0.045 & 0.108 \\
\hline
\multirow{5}{*}{\rotatebox{90}{Exchange}} & 12.5\% & \boldres{0.004} & \boldres{0.038} & 0.043 & 0.151 & 0.040 & 0.137 & 0.023 & 0.106 & 0.112 & 0.224 & 0.235 & 0.407 & 0.345 & 0.483 & 0.280 & 0.411 & \secondres{0.017} & \secondres{0.090} & 0.025 & 0.112 & 0.019 & 0.098 \\
& 25\% & \boldres{0.004} & \boldres{0.039} & 0.104 & 0.240 & 0.043 & 0.143 & 0.037 & 0.134 & 0.184 & 0.288 & 0.412 & 0.551 & 0.443 & 0.561 & 0.253 & 0.383 & \secondres{0.030} & \secondres{0.122} & 0.037 & 0.136 & 0.034 & 0.130 \\
& 37.5\% & \boldres{0.005} & \boldres{0.042} & 0.149 & 0.296 & 0.045 & 0.148 & 0.053 & 0.159 & 0.281 & 0.355 & 0.447 & 0.579 & 0.393 & 0.522 & 0.316 & 0.472 & \secondres{0.045} & \secondres{0.150} & 0.050 & 0.159 & 0.046 & 0.153 \\
& 50\% & \boldres{0.005} & \boldres{0.046} & 0.215 & 0.363 & \secondres{0.039} & 0.140 & 0.071 & 0.185 & 0.387 & 0.431 & 0.499 & 0.593 & 0.575 & 0.649 & 0.358 & 0.489 & 0.036 & \secondres{0.130} & 0.061 & 0.176 & 0.062 & 0.176 \\
& Avg & \boldres{0.005} & \boldres{0.041} & 0.128 & 0.262 & 0.042 & 0.142 & 0.046 & 0.146 & 0.241 & 0.324 & 0.398 & 0.532 & 0.439 & 0.554 & 0.302 & 0.439 & \secondres{0.032} & \secondres{0.123} & 0.043 & 0.146 & 0.040 & 0.139 \\
\hline
\multicolumn{2}{c|}{1st Count} & 25 & 32 & 0 & 0 & 0 & 0 & 0 & 0 & 3 & 1 & 0 & 0 & 0 & 0 & 7 & 2 & 0 & 0 & 0 & 0 & 0 & 0 \\
\hline
\end{tabular}
}
\end{table}

Table~\ref{tab:full_imputation_results} presents imputation results under different missing rates $r \in \{12.5\%, 25\%, 37.5\%, 50\%\}$. Following the experimental protocol in TimesNet~\citep{wu2023timesnet}, we randomly mask values in the test sequences and evaluate reconstruction quality using MSE and MAE metrics.

\subsection{Anomaly Detection}
\label{appx:anomaly_results}
\begin{table*}[h]
\centering
\small
\caption{Full results for the anomaly detection task. The P, R and F1 represent the precision, recall and F1-score (\%) respectively. F1-score is the harmonic mean of precision and recall. A higher value of P, R and F1 indicates a better performance.}
\label{tab:full_anomaly_detection_results}
\resizebox{\textwidth}{!}{
\begin{tabular}{c|ccc|ccc|ccc|ccc|ccc|c}
\toprule
{Models} &
\multicolumn{3}{c|}{\textbf{SMD}} &
\multicolumn{3}{c|}{\textbf{MSL}} &
\multicolumn{3}{c|}{\textbf{SMAP}} &
\multicolumn{3}{c|}{\textbf{SWaT}} &
\multicolumn{3}{c|}{\textbf{PSM}} &
\multirow{2}{*}{\textbf{Avg F1}} \\

\cmidrule(lr){2-4}
\cmidrule(lr){5-7}
\cmidrule(lr){8-10}
\cmidrule(lr){11-13}
\cmidrule(lr){14-16}
Metric & P & R & F1 & P & R & F1 & P & R & F1 & P & R & F1 & P & R & F1 & \\
\midrule
\scalebox{0.9}{LSTM (1997)} \scalebox{0.8}{\citeyear{hochreiter1997long}} & 78.52 & 65.47 & 71.41 & 78.04 & 86.22 & 81.93 & 91.06 & 57.49 & 70.48 & 78.06 & 91.72 & 84.34 & 69.24 & 99.53 & 81.67 & 77.97 \\
\scalebox{0.9}{Transformer} \scalebox{0.8}{\citeyear{vaswani2017attention}} & 83.58 & 76.13 & 79.56 & 71.57 & 87.37 & 78.68 & 89.37 & 57.12 & 69.70 & 68.84 & 96.53 & 80.37 & 62.75 & 96.56 & 76.07 & 76.88 \\
\scalebox{0.9}{LogTrans} \scalebox{0.8}{\citeyear{li2019enhancing}} & 83.46 & 70.13 & 76.21 & 73.05 & 87.37 & 79.57 & 89.15 & 57.59 & 69.97 & 68.67 & 97.32 & 80.52 & 63.06 & 98.00 & 76.74 & 76.60 \\
\scalebox{0.9}{TCN} \scalebox{0.8}{\citeyear{Franceschi2019UnsupervisedSR}}& 84.06 & 79.07 & 81.49 & 75.11 & 82.44 & 78.60 & 86.90 & 59.23 & 70.45 & 76.59 & 95.71 & 85.09 & 54.59 & 99.77 & 70.57 & 77.24 \\
\scalebox{0.9}{Reformer} \scalebox{0.8}{\citeyear{kitaev2020reformer}} & 82.58 & 69.24 & 75.32 & 85.51 & 83.31 & 84.40 & 90.91 & 57.44 & 70.40 & 72.50 & 96.53 & 82.80 & 59.93 & 95.38 & 73.61 & 77.31 \\
\scalebox{0.9}{Informer} \scalebox{0.8}{\citeyear{zhou2021informer}} & 86.60 & 77.23 & 81.65 & 81.77 & 86.48 & 84.06 & 90.11 & 57.13 & 69.92 & 70.29 & 96.75 & 81.43 & 64.27 & 96.33 & 77.10 & 78.83 \\
\scalebox{0.9}{Autoformer} \scalebox{0.8}{\citeyear{wu2021autoformer}} 88.06 & 82.35 & 85.11 & 77.27 & 80.92 & 79.05 & 90.40 & 58.62 & 71.12 & 89.85 & 95.81 & 92.74 & 99.08 & 88.15 & 93.29 & 84.26 \\
\scalebox{0.9}{Pyraformer} \scalebox{0.8}{\citeyear{liu2022pyraformer}} & 85.61 & 80.61 & 83.04 & 83.81 & 85.93 & 84.86 & 92.54 & 57.71 & 71.09 & 87.92 & 96.00 & 91.78 & 71.67 & 96.02 & 82.08 & 82.57 \\
\scalebox{0.9}{Anomaly Trans.} \scalebox{0.8}{\citeyear{xu2021anomaly}} & 88.91 & 82.23 & 85.49 & 79.61 & 87.37 & 83.31 & 91.85 & 58.11 & 71.18 & 72.51 & 97.32 & 83.10 & 68.35 & 94.72 & 79.40 & 80.50 \\
\scalebox{0.9}{Stationary} \scalebox{0.8}{\citeyear{liu2022non}} & 88.33 & 81.21 & 84.62 & 68.55 & 89.14 & 77.50 & 89.37 & 59.02 & 71.09 & 68.03 & 96.75 & 79.88 & 97.82 & 96.76 & 97.29 & 82.08 \\
\scalebox{0.9}{LSSL} \scalebox{0.8}{\citeyear{gu2022efficiently}} & 78.51 & 65.32 & 71.31 & 77.55 & 88.18 & 82.53 & 89.43 & 53.43 & 66.90 & 79.05 & 93.72 & 85.76 & 66.02 & 92.93 & 77.20 & 76.74 \\
\scalebox{0.9}{ETSformer} \scalebox{0.8}{\citeyear{woo2022etsformer}} & 87.44 & 79.23 & 83.13 & 85.13 & 84.93 & 85.03 & 92.25 & 55.75 & 69.50 & 90.02 & 80.36 & 84.91 & 99.31 & 85.28 & 91.76 & 82.87 \\
\scalebox{0.9}{DLinear} \scalebox{0.8}{\citeyear{zeng2023dlinear}} & 83.62 & 71.52 & 77.10 & 84.34 & 85.42 & 84.88 & 92.32 & 55.41 & 69.26 & 80.91 & 95.30 & 87.52 & 98.28 & 89.26 & 93.55 & 82.46 \\
\scalebox{0.9}{TiDE} \scalebox{0.8}{\citeyear{das2023long}} & 76.00 & 63.00 & 68.91 & 84.00 & 60.00 & 70.18 & 88.00 & 50.00 & 64.00 & 98.00 & 63.00 & 76.73 & 93.00 & 92.00 & 92.50 & 74.46 \\
\scalebox{0.9}{PatchTST} \scalebox{0.8}{\citeyear{nie2023patchtst}} & 87.26 & 82.14 & 84.62 & 88.34 & 70.96 & 78.70 & 90.64 & 55.46 & 68.82 & 91.10 & 80.94 & 85.72 & 98.84 & 93.47 & 96.08 & 82.79 \\
\scalebox{0.9}{iTransformer} \scalebox{0.8}{\citeyear{liu2023itransformer}} & 78.45 & 65.10 & 71.15 & 86.15 & 62.65 & 72.54 & 90.67 & 52.96 & 66.87 & 99.96 & 65.55 & 79.18 & 95.65 & 94.69 & 95.17 & 76.98 \\
\midrule
\scalebox{0.9}{\textbf{MAS4TS (Ours)}} & 86.76 & 83.51 & 85.10 & 90.65 & 74.96 & 82.06 & 89.87 & 55.44 & 68.58 & 91.61 & 94.74 & 93.14 & 91.87 & 98.44 & 95.04 & \textbf{84.78} \\
\bottomrule
\end{tabular}
}
\end{table*}

Table~\ref{tab:full_anomaly_detection_results} reports anomaly detection results on five benchmark datasets using precision (P), recall (R), and F1-score (F1) metrics. Following standard practice~\citep{xu2021anomaly}, we apply the point-adjustment protocol where a detection is considered correct if any point within an anomalous segment is identified.